\documentclass{article}

\usepackage[utf8]{inputenc}
\usepackage{algpseudocode}
\usepackage{algorithm}

\title{IFT - Pseudocode}
\author{Lutz; Paula}

\begin{document}
    
\maketitle

\section*{Statistical Analysis - Pseudocode}

%%% Main function for managing the statistical analysis using IFT
\begin{algorithm}
    \caption{StatAnalysis (Statistical Analysis using IFT-functions in SA class)}
    \begin{algorithmic}[1]
        \Require Netlist; OG Verilog + \_tb
        \Ensure external simulation runs through
        \State IFTGenIFTVerilog(verilogPath) \Comment{generates IFT extended Verilog}
        \State IFTPathAnalysis()
    \end{algorithmic}
\end{algorithm}
%%%%%%%%%%%%%%%%%%%%%%%%%%%%%%%%%%%%%%%%%%%%%%%%%%%%%%%%%%%%%%%%

%%% Generation functions for exporting the IFT extended Verilog and testbench
\begin{algorithm}
    \caption{IFTGenVerilog (IFTGeneration of extended Verilog)}
    \begin{algorithmic}[1]
        \Require OG Verilog + \_tb
        \State design.exportIFTVerilog()
            \For{m in Modules}
                \State exportIFTVerilogModuleHeader()
                \State exportIFTVerilogModuleInternals()
            \EndFor
    \end{algorithmic}
\end{algorithm}

\begin{algorithm}
    \caption{IFTGenTestBench (IFTGeneration of IFT testbench)}
    \begin{algorithmic}[1]
        \Require OG Verilog + \_tb
        \State generateInputVCD() \Comment{simulates OG Verilog + \_tb}
        \State exportIFTTestBench(InputVCD)
    \end{algorithmic}
\end{algorithm}
%%%%%%%%%%%%%%%%%%%%%%%%%%%%%%%%%%%%%%%%%%%%%%%%%%%%%%%%%%%%%%%%

%%% Executes the simulation of the extended IFT verilog design
\begin{algorithm}
    \caption{IFTExe (IFTExecution)}
    \begin{algorithmic}[1]
        \Require IFT Verilog
        \State SimulateIFTVerilog(IFTVerilogPath + TestBenchPath)
        \State readBackVCD(outputVCD) \Comment{stores simulation values in Design}
    \end{algorithmic}
\end{algorithm}
%%%%%%%%%%%%%%%%%%%%%%%%%%%%%%%%%%%%%%%%%%%%%%%%%%%%%%%%%%%%%%%%

%%%
%%% Calls TagGen() which decides which inputs get tagged and writes the tag input file
%%% Calls testbench creation to rewrite testbench variable
\begin{algorithm}
    \caption{IFTPA (IFT Path Analysis)}
    \begin{algorithmic}[1]
        \Require OR IFT Verilog + \_tb
        \State TagGen() \Comment{generate TAG inputs}
        \For{i in 100}
            \State IFTGenTestbench() \Comment{At some point change of design input values}
            \State IFTExe(TAGs, TestBenchPath)
            \State StatisticalAnalysis \Comment identify + enumerate paths from simulation values
        \EndFor 
    \end{algorithmic}
\end{algorithm}
%%%%%%%%%%%%%%%%%%%%%%%%%%%%%%%%%%%%%%%%%%%%%%%%%%%%%%%%%%%%%%%%

%%% a method to decide which inputs of the current design get a tag for the next simulation run
%%% either generate a new set of tags every run or prepare the full set of tags and inputs beforehand
\begin{algorithm}
    \caption{TagGen (Tag Generation method)}
    \begin{algorithmic}[1]
        \Require Number of Inputs/Tags in design
        \State DecideTags() \Comment some method needed to choose which tags to set
            \State - Random (numberOfTagsIn\%) \Comment{Choose Tags at random (predetermined number of tags)}
            \State - Analyze InputNames (special use cases (designs/design groups)) \Comment{decide which inputs to tag on NameDatabase??}
            \State - User Selection
            \State - PreProcessing of Variable Paths \Comment{Choose which variables to tag based on criteria for their paths}
    \end{algorithmic}
\end{algorithm}
%%%%%%%%%%%%%%%%%%%%%%%%%%%%%%%%%%%%%%%%%%%%%%%%%%%%%%%%%%%%%%%%
%%%
%%% Calls TagGen() which decides which inputs get tagged and writes the tag input file
%%% Calls testbench creation to rewrite testbench variable
\begin{algorithm}
    \caption{IFTPA (IFT Path Analysis)}
    \begin{algorithmic}[1]
        \Require OR IFT Verilog + \_tb
        \State TagGen() \Comment{generate TAG inputs}
        \For{i in numSimulationRuns}
            \State IFTGenTestbench() \Comment{At some point change of design input values}
            \State IFTExe(TAGs, TestBenchPath)
            \State StatisticalAnalysis \Comment identify + enumerate paths from simulation values
        \EndFor 
    \end{algorithmic}
\end{algorithm}
%%%%%%%%%%%%%%%%%%%%%%%%%%%%%%%%%%%%%%%%%%%%%%%%%%%%%%%%%%%%%%%%

\newpage

\section*{Parameters}

For Information flow tracking with subsequent statistical analysis of the obtained results a lot of possible parameters exist.\\
How are these parameters set? automatically or manually?\\

\paragraph*{Parameters directly relevant for the tags}
\begin{itemize}
    \item number of tags
    \item length of time tags are set (how long?)
    \item at which point in time are tags set (when?)
\end{itemize}

\paragraph*{Parameters for the statistical analysis}
\begin{itemize}
    \item number of simulation runs
    \item number of input variable combinations 
    \item input and tag combinations
\end{itemize}

\subsubsection*{input and tag combinations}
How are input and tag combinations decided?\\
Are multiple new tags set for the same inputs?\\
How many sets of inputs are needed for meaningful stastitical analysis?\\

\end{document}